The Alabama Paradox --- in which increasing the House size causes a
state to \emph{lose} a seat --- was documented in the 1880 Census
and is the most politically embarrassing property of the Hamilton
(largest remainders) apportionment method.
Two related paradoxes complete the triad of classical apportionment
anomalies: the Population Paradox (a faster-growing state loses a seat
to a slower-growing state) and the New-States Paradox (adding a state
with proportional seats causes existing states to exchange seats).
All three paradoxes afflict Hamilton but not the divisor methods
(Huntington-Hill, Webster, Adams, Jefferson).

This paper proves, with formal mathematics and numerical examples,
(1) that Hamilton is susceptible to all three paradoxes,
(2) that all four divisor methods are immune to all three paradoxes,
and (3) that no apportionment method can be simultaneously
quota-compliant and paradox-immune --- the Balinski-Young impossibility
theorem (1982).

For Hamilton susceptibility we give numerical examples:
the 1880 Alabama Paradox (adding a seat reduced Alabama from 8 to 7),
a 1900 Population Paradox example from U.S.\ census history,
and the 1907 Oklahoma New-States Paradox.
For divisor method immunity we give a formal proof using the
monotone priority-queue structure: in a priority-queue divisor method,
adding a seat or increasing a state's population can only add seats ---
no existing seat allocation is reconsidered.
The Balinski-Young theorem is proved by showing that satisfying the
quota rule for all possible population vectors forces susceptibility
to the Alabama paradox: the fractional-remainder ranking that ensures
quota compliance is necessarily non-monotone in House size.
The paper concludes with the historical and policy consequences of
these mathematical results for Congress's 1941 adoption of
Huntington-Hill.
